%% paper_expanded.tex
%% Expanded version of the LGDS short note
%% Title: Common Optimal Finite-Rank Projections: Coherence, Coverability, and Regret
%% Format: elsarticle (compatible with Systems & Control Letters / Automatica)

\documentclass[final,3p,times]{elsarticle}

\usepackage{amsmath,amssymb,amsthm}
\usepackage{booktabs}
\usepackage{microtype}
\usepackage[colorlinks=true,citecolor=blue,urlcolor=blue,linkcolor=black]{hyperref}

% Auto-generated by lgds.py -- do not edit by hand

\newcommand{\CondAAngleOne}{54.1}
\newcommand{\CondAAngleTwo}{24.0}
\newcommand{\CondARegretUoneTa}{0.000}
\newcommand{\CondARegretUoneTb}{0.036}
\newcommand{\CondARegretUtwoTa}{0.763}
\newcommand{\CondARegretUtwoTb}{0.000}
\newcommand{\CondARegretUjointTa}{0.026}
\newcommand{\CondARegretUjointTb}{0.026}
\newcommand{\GammaCondA}{0.026}
\newcommand{\CondBMaxAngleDeg}{0.000}
\newcommand{\CondBPertEps}{0.1}
\newcommand{\CondBPertDraws}{20}
\newcommand{\CondBPertMaxRad}{0.102}
\newcommand{\CondBPertThreshold}{0.15}
\newcommand{\HiDimFRandMean}{0.118}
\newcommand{\HiDimFPB}{0.617}
\newcommand{\HiDimFPF}{0.619}
\newcommand{\HiDimRegret}{0.0023}
\newcommand{\HiDimCoarse}{0.362}
\newcommand{\HiDimInstDep}{0.144}
\newcommand{\HiDimUpper}{1.015}


\newtheorem{theorem}{Theorem}
\newtheorem{proposition}[theorem]{Proposition}
\newtheorem{corollary}[theorem]{Corollary}
\newtheorem{lemma}[theorem]{Lemma}
\newtheorem{definition}[theorem]{Definition}
\newtheorem{remark}[theorem]{Remark}

\newcommand{\Gr}{\mathrm{Gr}}
\newcommand{\tr}{\mathrm{tr}}
\newcommand{\R}{\mathbb{R}}
\newcommand{\E}{\mathbb{E}}
\newcommand{\rank}{\mathrm{rank}}
\newcommand{\ran}{\mathrm{ran}}
\newcommand{\Reg}{\mathrm{Reg}}
\newcommand{\Emax}{E_{\max}}
\renewcommand{\qedsymbol}{$\blacksquare$}

\journal{Systems \& Control Letters}

\begin{document}

\begin{frontmatter}

\title{Common Optimal Finite-Rank Projections:\\
       Coherence, Coverability, and Regret}

\author{Kunal Bhatia}
\address{Independent Researcher, Heidelberg, Germany\\
ORCID: \href{https://orcid.org/0009-0007-4447-6325}{0009-0007-4447-6325}}

\begin{abstract}
For a finite family of positive semidefinite (PSD) target operators on
$\R^d$ and a fixed rank $r$, we ask when a single rank-$r$ orthogonal
projection is simultaneously optimal for every target under two natural
scoring rules: the \emph{predictive} Ky--Fan trace
$S^{\mathrm{pred}}_t(P)=\tr(P M_t)$ and the \emph{causal} operator norm
$S^{\mathrm{caus}}_t(P)=\lambda_{\max}(P K_t P)$.
Existence of a common optimal projection under $S^{\mathrm{pred}}$ is
equivalent to \emph{$r$-coherence}---all $M_t$ share a common dominant
$r$-dimensional invariant subspace.
Existence of a common optimal projection under $S^{\mathrm{caus}}$ is
equivalent to \emph{$r$-coverability}---a single rank-$r$ subspace
nontrivially intersects the leading eigenspace of every $K_t$.
$r$-coherence implies $r$-coverability but the converse fails; a $3\times 3$
witness exhibits the strict gap.
When coverability fails over a regime of local intervention Jacobians
$\{J_i\}$, the eigenvalue-weighted surrogate
$B_\lambda=\sum_i w_i\lambda_1(J_iJ_i^\top)u_iu_i^\top$ provides a
Ky--Fan-tractable approximation of the causal landscape; we prove a
sandwich
$\tilde F(P)\le F(P)\le\tilde F(P)+\sum_i w_i\lambda_2(J_iJ_i^\top)$
and an instance-dependent regret bound
$F(P_F^\star)-F(P_B^\star)\le\sum_i w_i\lambda_2(J_iJ_i^\top)(1-\|P_F^\star u_i\|^2)$
that is tight in the rank-$1$, strong-glue, and common-shape regimes.
Numerical illustrations include a stationary linear Gaussian dynamical
system certifying non-coherence ($\Gamma=\GammaCondA>0$), the $3\times 3$
fork witness in which a common optimal projection exists under
$S^{\mathrm{caus}}$ but not under $S^{\mathrm{pred}}$, and a $3\times 3$
weak-glue counterexample in which the coarse regret bound overestimates
the true regret by a factor of $36$.
\end{abstract}

\begin{keyword}
finite-rank projection \sep Grassmannian \sep Ky Fan theorem \sep
$r$-coherence \sep $r$-coverability \sep operator-norm objective
\sep regret bound
\MSC[2020] 93B07 \sep 15A18 \sep 93E10 \sep 90C26
\end{keyword}

\end{frontmatter}

% =========================================================================
\section{Introduction}
% =========================================================================

A recurrent structural question in observer design, sensor selection, and
reduced-order modelling is whether a single low-rank linear projection can
serve simultaneously as the optimal observation map for an entire family of
prediction or intervention tasks, or whether the optimal projection is
inherently task-dependent.
The answer turns out to depend not only on the family of targets but on
\emph{which scoring rule} is used to measure observer quality.

For a fixed projector rank $r$ and a state space $\R^d$, two scoring rules
are natural and arise in different contexts:
\begin{itemize}
\item the \emph{predictive trace} $S^{\mathrm{pred}}_t(P) = \tr(P M_t)$,
which corresponds to squared-error prediction risk in linear Gaussian
settings and is governed by the Ky Fan variational principle
\cite{Fan1949};
\item the \emph{causal operator norm} $S^{\mathrm{caus}}_t(P) =
\lambda_{\max}(P K_t P)$, which corresponds to the strongest accessible
intervention response under projector $P$, with $K_t$ a local
output Gram or causal-response operator.
\end{itemize}

This paper establishes, for arbitrary finite PSD families $\{M_t\}$ and
$\{K_t\}$, a structural answer to the existence question for each
scoring rule, and characterises the gap between them.

\paragraph{Contribution}
The contributions are four:

\begin{enumerate}
\item \emph{Predictive simultaneous optimality} (Theorem~\ref{thm:pred}):
A single rank-$r$ projection is simultaneously optimal for all predictive
targets $\{M_t\}$ if and only if the family is \emph{$r$-coherent}: all
$M_t$ share a common dominant $r$-dimensional invariant subspace
(equivalently, share a top-$r$ eigenspace under strict $r$-gap).
This is a direct application of Ky Fan trace maximisation; the LGDS
Bayes-risk instance is recovered as a corollary.

\item \emph{Causal simultaneous optimality} (Theorem~\ref{thm:caus}):
A single rank-$r$ projection is simultaneously optimal for all causal
targets $\{K_t\}$ under the operator-norm score if and only if the
family is \emph{$r$-coverable}: some rank-$r$ subspace intersects
$\Emax(K_t)$ nontrivially for every $t$.
The proof uses the rank-$1$-driven structure of the
$\lambda_{\max}(P K P)$ maximiser.

\item \emph{The fork} (Theorem~\ref{thm:fork}): For a single PSD family,
$r$-coherence implies $r$-coverability strictly.
A $3\times 3$ witness shows the gap: a single common optimal projection
can exist under the operator-norm score while no common optimal
projection exists under the trace score.
The dissociation is algebraic, not a consequence of approximation error.

\item \emph{Weak-glue regret} (Theorems~\ref{thm:weakglue},
\ref{thm:instancereg}): When coverability fails, the
eigenvalue-weighted surrogate
$B_\lambda = \sum_i w_i \lambda_1(A_i) u_i u_i^\top$ converts the
non-Ky--Fan operator-norm objective into a Ky--Fan eigenproblem.
The surrogate is sandwiched between $F$ and $F + \sum_i w_i\lambda_2(A_i)$,
giving a coarse regret bound; the instance-dependent refinement
$\sum_i w_i \lambda_2(A_i)(1-\|P_F^\star u_i\|^2)$ unifies the
equality regimes (rank-$1$ Jacobians, strong glue, common-shape) and
is provably tighter than the coarse bound.
\end{enumerate}

\paragraph{Position in the literature}
The Grassmannian $\Gr(r,d)$ is the standard parameter space for rank-$r$
linear maps, with smooth Riemannian structure inherited from $O(d)$
\cite{Edelman1998}; principal-subspace optimality and the Ky Fan trace
maximisation are classical \cite{Fan1949,HornJohnson2013,Hotelling1933}.
The condition that several covariance-like matrices share a common
eigenstructure has a long history under the names \emph{common principal
components} \cite{Flury1984,Krzanowski1979} and \emph{joint
diagonalisation} \cite{Pham2001}.
Reduced-order models for linear systems are commonly built from balanced
Gramians \cite{Moore1981} or by truncating in $H_2$/$H_\infty$ norms
\cite{Antoulas2005}; sensor-placement and output-design methods optimise
estimation or reconstruction quality for a single objective
\cite{JoshiBoyd2009}.
Robust subspace recovery \cite{LermanMaunu2018} addresses outlier and
contamination structure, complementary to the exact-optimality focus
here.
Within this landscape, Theorems~\ref{thm:pred} and~\ref{thm:caus} are
structural consequences of the Ky Fan trace and operator-norm variational
principles respectively; the contribution of this paper is to put the
two existence laws side by side, prove the strict gap between them
(Theorem~\ref{thm:fork}), and develop the weak-glue regret comparison
(Theorems~\ref{thm:weakglue}, \ref{thm:instancereg}) that quantifies
suboptimality when neither exact condition holds.

\paragraph{Outline}
Section~\ref{sec:setup} gives the setup, the two scoring rules, and the
notion of a common optimal rank-$r$ projection.
Section~\ref{sec:pred} proves the predictive theorem and instantiates it
in LGDS, recovering the Bayes-risk formula and the regret-angle bound.
Section~\ref{sec:caus} proves the causal theorem.
Section~\ref{sec:fork} proves the fork.
Section~\ref{sec:weakglue} develops the weak-glue surrogate.
Section~\ref{sec:numerics} gives numerical illustrations.
Section~\ref{sec:scope} closes with scope and discussion.

% =========================================================================
\section{Setup}
\label{sec:setup}
% =========================================================================

\paragraph{State space and projector class}
The state space is $\R^d$ with the Euclidean inner product.
A \emph{rank-$r$ projection} ($1 \le r < d$) is an orthogonal projector
$P$ of rank $r$; equivalently, $P = U^\top U$ for a row-orthonormal
$U \in \R^{r\times d}$ with $UU^\top = I_r$.
The set of rank-$r$ projectors is naturally identified with the
Grassmannian $\Gr(r,d)$, a smooth compact Riemannian manifold of
dimension $r(d-r)$ with metric and geodesic structure inherited from
$O(d)$ \cite{Edelman1998}.
Two projectors with the same range are equivalent on $\Gr(r,d)$.

\paragraph{Target family and scores}
A \emph{target family} $\mathcal{T}$ indexes the tasks.
Each $t\in\mathcal{T}$ induces two PSD operators on $\R^d$:
\begin{itemize}
\item a \emph{predictive target operator} $M_t \succeq 0$,
\item a \emph{causal target operator} $K_t \succeq 0$.
\end{itemize}
The two scoring rules considered are:
\begin{equation}\label{eq:scores}
S^{\mathrm{pred}}_t(P) := \tr(P M_t),
\qquad
S^{\mathrm{caus}}_t(P) := \lambda_{\max}(P K_t P).
\end{equation}
The predictive score is the Ky Fan trace; it is a smooth (analytic)
function of $P\in\Gr(r,d)$.
The causal score is the operator norm of $P K_t P$; it is continuous and
$1$-Lipschitz in the operator norm of $P$, but smooth only on the
Grassmannian stratum where the leading eigenvalue of $P K_t P$ is simple,
with non-smooth points at eigenvalue crossings \cite{HornJohnson2013}.

\paragraph{Spectral subspaces}
\begin{definition}[Dominant invariant subspace; leading eigenspace]
\label{def:domsubspace}
For PSD $T\succeq 0$ on $\R^d$ with eigenvalues
$\lambda_1(T)\ge\cdots\ge\lambda_d(T)\ge 0$, a \emph{dominant
$r$-dimensional invariant subspace} of $T$ is any $r$-dimensional
subspace $V\subseteq\R^d$ satisfying
$\tr(\Pi_V T)=\sum_{i=1}^r\lambda_i(T)$ (the Ky Fan maximum).
Under strict $r$-gap $\lambda_r(T)>\lambda_{r+1}(T)$, $V$ is unique and
equals the span of eigenvectors of $\lambda_1(T),\ldots,\lambda_r(T)$.
The \emph{leading eigenspace} is
$\Emax(T):=\ker(T-\lambda_1(T)I)$ when $T\ne 0$, and
$\Emax(0):=\R^d$ by convention.
\end{definition}

\paragraph{Common optimal projection}
\begin{definition}[Common optimal rank-$r$ projection]\label{def:fgp}
A rank-$r$ orthogonal projector $P^\star$ is a \emph{common optimal
rank-$r$ projection} for the family $\mathcal{T}$ under score $S$ if
\begin{equation}
  S_t(P^\star) \;=\; \max_{P:\,P^\top=P=P^2,\,\rank P=r}\,S_t(P)
  \qquad\text{for every }t\in\mathcal{T}.
\end{equation}
Throughout this paper, $P$ ranges over rank-$r$ orthogonal projectors
(equivalently $P\in\Gr(r,d)$ via $P=U^\top U$, $UU^\top=I_r$); the
$\max$ is over this set.
The same property is sometimes called \emph{finite global privilege} of
$P^\star$ for $\mathcal{T}$ under $S$; we use the two terms
interchangeably.
\end{definition}

The central question of the paper is when a common optimal rank-$r$
projection exists, under each scoring rule, and what structure is forced
on the family if so.
The two scoring rules give different answers.

% =========================================================================
\section{Predictive Global Privilege via \texorpdfstring{$r$}{r}-coherence}
\label{sec:pred}
% =========================================================================

\subsection{The general PSD theorem}

\begin{definition}[$r$-coherence]\label{def:rcoh}
A PSD family $\{M_t\}_{t\in\mathcal{T}}$ on $\R^d$ is \emph{$r$-coherent}
if there exists a rank-$r$ subspace $V^\star\subseteq\R^d$ contained
in a dominant $r$-dimensional invariant subspace of every $M_t$.
Under strict $r$-gaps ($\lambda_r(M_t)>\lambda_{r+1}(M_t)$ for all $t$),
this is equivalent to the equality of the top-$r$ eigenspaces of all
$M_t$.
\end{definition}

\begin{theorem}[Predictive simultaneous optimality]\label{thm:pred}
A common optimal rank-$r$ projection for the PSD family $\{M_t\}$ under
$S^{\mathrm{pred}}$ exists if and only if $\{M_t\}$ is $r$-coherent.
\end{theorem}

\begin{proof}
By the Ky Fan variational principle \cite{Fan1949,HornJohnson2013}, for
each $t$,
\begin{equation}\label{eq:kyfan}
\max_{\rank P = r} \tr(P M_t) = \sum_{i=1}^r \lambda_i(M_t),
\end{equation}
attained exactly by projectors whose range lies within a dominant
$r$-dimensional invariant subspace of $M_t$.
($\Leftarrow$) If $V^\star$ is a common rank-$r$ subspace contained in
every dominant invariant subspace, then $P^\star := \Pi_{V^\star}$ attains
the maximum~\eqref{eq:kyfan} for every $t$.
($\Rightarrow$) Conversely, if $P^\star$ attains the maximum for every
$t$, then $\ran(P^\star)$ lies within a dominant $r$-dimensional invariant
subspace of every $M_t$, so the family is $r$-coherent.
Under strict $r$-gaps the dominant invariant subspace is unique and equal
to the top-$r$ eigenspace, so all $M_t$ share the same top-$r$ eigenspace.
\end{proof}

\begin{remark}[Strictness of $r$-coherence]
$r$-coherence is strictly weaker than simultaneous diagonalisability.
Pairwise commutativity $[M_s,M_t]=0$ implies a common eigenbasis and hence
$r$-coherence for every $r$, but $r$-coherence with a fixed $r$ does not
require commutativity: only the top-$r$ eigenspace must coincide.
\end{remark}

\begin{remark}[Genericity]
For generic $\{M_t\}$, distinct $M_t$ have distinct top-$r$ eigenspaces.
$r$-coherence is therefore a non-generic algebraic condition; outside
$r$-coherence, the predictive optimum is target-dependent.
\end{remark}

\subsection{LGDS instantiation}

The predictive theorem's classical instance is the stationary linear
Gaussian dynamical system (LGDS) with squared-error prediction loss.
Let
\begin{equation}
x_{t+1} = A x_t + \xi_t, \qquad \xi_t\sim\mathcal{N}(0,Q),
\end{equation}
with $A$ stable and $Q\succ 0$, and let $\Sigma\succ 0$ denote the
stationary covariance solving $A\Sigma A^\top+Q=\Sigma$ \cite{Kailath1980}.
A prediction target is a pair $(C,\tau)$ with $C\in\R^{p\times d}$ and
horizon $\tau\ge 1$; the goal is to predict $Cx_{t+\tau}$ from a rank-$r$
projection of the current state.
The map $P=U^\top U$ is an orthogonal projection in the
\emph{whitened} coordinates $\Sigma^{-1/2}x$; in original state
coordinates the observed subspace is its $\Sigma^{-1/2}$-pullback,
and the observation map is $\phi_U(x)=U\Sigma^{-1/2}x$.
Pre-whitening places all candidate projections on a common isotropic
footing.

\begin{lemma}[Bayes risk for LGDS]\label{lem:bayes}
For projector $P=U^\top U$ with $UU^\top=I_r$, define the linear map
$\phi_U(x)=U\Sigma^{-1/2}x$.
At stationarity, the squared-error Bayes risk for predicting $Cx_{t+\tau}$
from $\phi_U(x_t)$ is
\begin{equation}\label{eq:lgdsbayes}
\mathcal{R}(U;\,C,\tau)
= \tr(C\Sigma C^\top) - \tr\!\bigl(U M_{C,\tau} U^\top\bigr),
\end{equation}
where the \emph{target information matrix} is
\begin{equation}\label{eq:M}
M_{C,\tau} = \Sigma^{1/2}(A^\tau)^\top C^\top C A^\tau \Sigma^{1/2}.
\end{equation}
\end{lemma}

\begin{proof}
$x_{t+\tau}=A^\tau x_t+w_\tau$ with $w_\tau$ independent of $x_t$, so
$\E[x_{t+\tau}x_t^\top]=A^\tau\Sigma$.
With $z=U\Sigma^{-1/2}x_t$, $\mathrm{Cov}(z,z^\top)=I_r$ and
$\mathrm{Cov}(Cx_{t+\tau},z^\top)=CA^\tau\Sigma^{1/2}U^\top$.
The Gaussian linear MMSE formula
$\mathcal{R}=\tr(\mathrm{Var}(Cx_{t+\tau}))-\tr(\mathrm{Cov}(Cx_{t+\tau},z^\top)\mathrm{Cov}(z,z^\top)^{-1}\mathrm{Cov}(z,Cx_{t+\tau}))$
collapses to~\eqref{eq:lgdsbayes}.
\end{proof}

Since $\tr(C\Sigma C^\top)$ is independent of $U$, minimising
$\mathcal{R}(U;C,\tau)$ is equivalent to maximising
$\tr(U M_{C,\tau} U^\top)=\tr(P M_{C,\tau})$ where $P=U^\top U$.
Theorem~\ref{thm:pred} applied to the LGDS family
$\{M_{C_k,\tau_k}\}$ gives:

\begin{corollary}[LGDS predictive iff]
A common rank-$r$ projection minimising squared-error prediction risk for
every $(C_k,\tau_k)\in\mathcal{T}$ exists if and only if
$\{M_{C_k,\tau_k}\}$ is $r$-coherent.
\end{corollary}

\subsection{Task regret and the angle bound}

The deviation of $P$ from $P^\star$ admits a sharp two-sided
characterisation in principal angles \cite{BjorckGolub1973}.
Let $M$ be PSD with eigenvalues $\lambda_1\ge\cdots\ge\lambda_d$ and
top-$r$ eigenspace projector $P^\star$.
Define the \emph{task regret}
\begin{equation}\label{eq:regret}
\Reg(P;\,M) := \tr(P^\star M) - \tr(P M)
= \sum_{i=1}^r\lambda_i(M)-\tr(P M)\;\ge\;0.
\end{equation}

\begin{corollary}[Regret--angle bound]\label{cor:regretbound}
Let $\|\sin\Theta\|_F^2 := \|\sin\Theta(P,P^\star)\|_F^2$ be the squared
Frobenius norm of the principal-angle sines between $\ran(P)$ and
$\ran(P^\star)$.
Then
\begin{equation}\label{eq:regretbound}
(\lambda_r-\lambda_{r+1})\|\sin\Theta\|_F^2
\;\le\;\Reg(P;\,M)\;\le\;
(\lambda_1-\lambda_d)\|\sin\Theta\|_F^2.
\end{equation}
\end{corollary}

\begin{proof}
Let $v_i$ be unit eigenvectors of $M$ with eigenvalues $\lambda_i$, and
$u_k$ orthonormal vectors spanning $\ran(P)$.
Set $p_i:=\sum_{k=1}^r(v_i^\top u_k)^2$.
Then $\tr(PM)=\sum_i\lambda_i p_i$ and $\sum_i p_i=r$, while
$\|\sin\Theta\|_F^2=\sum_{i\le r}(1-p_i)=\sum_{i>r}p_i$.
Centering at $\lambda_{r+1}$:
\[
\Reg(P;\,M)
=\sum_{i\le r}(\lambda_i-\lambda_{r+1})(1-p_i)
+\sum_{i>r}(\lambda_{r+1}-\lambda_i)\,p_i.
\]
Both sums are nonnegative.
The lower bound follows by dropping the second sum and bounding
$\lambda_i-\lambda_{r+1}\ge\lambda_r-\lambda_{r+1}$.
The upper bound follows by bounding the two coefficients by $\lambda_1
-\lambda_{r+1}$ and $\lambda_{r+1}-\lambda_d$ respectively.
\end{proof}

\begin{corollary}[Approximate predictive coherence]\label{cor:approxcoh}
For a finite predictive family $\{M_a\}$, define the \emph{minimax
predictive regret}
\begin{equation}\label{eq:gamma}
\Gamma_{\mathrm{pred}}
:= \min_{P\in\Gr(r,d)}\,\max_a\,\Reg(P;\,M_a).
\end{equation}
$\Gamma_{\mathrm{pred}}=0$ iff $\{M_a\}$ is $r$-coherent.
If there exists $P\in\Gr(r,d)$ with $\|\sin\Theta(P,P_a^\star)\|_F^2
\le\varepsilon^2$ for every $a$, then
$\Gamma_{\mathrm{pred}}\le\Lambda\,\varepsilon^2$, where
$\Lambda:=\max_a[\lambda_1(M_a)-\lambda_d(M_a)]$.
\end{corollary}

\begin{proof}
The first claim is Theorem~\ref{thm:pred} applied to the family.
The second applies the upper bound of Corollary~\ref{cor:regretbound}
uniformly in $a$.
\end{proof}

% =========================================================================
\section{Causal Global Privilege via \texorpdfstring{$r$}{r}-coverability}
\label{sec:caus}
% =========================================================================

The causal score $S^{\mathrm{caus}}_t(P)=\lambda_{\max}(P K_t P)$ is not
a Ky Fan trace; the variational principle that governs it is different.
The structural condition that emerges is correspondingly different.

\subsection{Local maximiser structure}

\begin{proposition}[Fixed-target causal maximiser]\label{prop:localmax}
For fixed PSD $K\succeq 0$ on $\R^d$ with leading invariant subspace
$\Emax(K)$ (the eigenspace of $\lambda_1(K)$),
\begin{equation}\label{eq:localmax}
\arg\max_{\rank P=r}\,\lambda_{\max}(P K P)
=\bigl\{P:\rank(P)=r,\;\ran(P)\cap\Emax(K)\ne\{0\}\bigr\},
\end{equation}
and the maximum value is $\lambda_1(K)$.
\end{proposition}

\begin{proof}
For any unit $v\in\R^d$, $v^\top K v\le\lambda_1(K)$, with equality iff
$v\in\Emax(K)$.
For rank-$r$ $P$,
$\lambda_{\max}(PKP)=\sup_{v\in\ran(P),\,\|v\|=1}v^\top K v$.
This supremum equals $\lambda_1(K)$ iff $\ran(P)$ contains a unit vector
in $\Emax(K)$, i.e., $\ran(P)\cap\Emax(K)\ne\{0\}$.
\end{proof}

The contrast with the predictive case is sharp: under
$S^{\mathrm{pred}}=\tr(PM)$ the optimal $P$ must have its range
\emph{equal to} a dominant $r$-dimensional invariant subspace of $M$;
under $S^{\mathrm{caus}}=\lambda_{\max}(PKP)$ it suffices that
$\ran(P)$ \emph{intersects} the leading invariant subspace of $K$
nontrivially.
The local problem is rank-$1$-driven regardless of $r$.

\subsection{\texorpdfstring{$r$}{r}-coverability and the causal theorem}

\begin{definition}[$r$-coverability]\label{def:cov}
A PSD family $\{K_t\}$ is \emph{$r$-coverable} if there exists a rank-$r$
subspace $V\subseteq\R^d$ with $V\cap\Emax(K_t)\ne\{0\}$ for every $t$.
When each $K_t$ has a simple top eigenvalue with top eigenvector $u_t$,
$r$-coverability reduces to $\dim\mathrm{span}\{u_t\}\le r$.
\end{definition}

\begin{theorem}[Causal simultaneous optimality]\label{thm:caus}
A common optimal rank-$r$ projection for the PSD family $\{K_t\}$ under
$S^{\mathrm{caus}}$ exists if and only if $\{K_t\}$ is $r$-coverable.
\end{theorem}

\begin{proof}
By Proposition~\ref{prop:localmax}, $P^\star$ is a common optimal rank-$r$
projection under $S^{\mathrm{caus}}$ iff $\ran(P^\star)\cap\Emax(K_t)\ne\{0\}$
for every $t$.
Existence of such a rank-$r$ subspace is the definition of
$r$-coverability.
\end{proof}

\begin{remark}[Coverability is much weaker than coherence]
$r$-coherence demands that one rank-$r$ subspace lie \emph{within} the
dominant invariant subspace of every $M_t$.
$r$-coverability demands only that the rank-$r$ subspace \emph{hit}
each $\Emax(K_t)$.
The latter is satisfied whenever the union of leading directions
$\{u_t\}$ spans at most $r$ dimensions.
\end{remark}

\subsection{Strong glue: a sufficient condition}

\begin{corollary}[Strong glue]\label{cor:strongglue}
Suppose $\{K_t\}_{t\in T}$ is indexed by a finite-support measure
$\mu=\sum_i w_i\delta_{x_i}$ on a state space, with $K_i:=K_{x_i}$ having
simple top eigenvalues and top eigenvectors $u_i$.
Then $\{K_i\}$ is $r$-coverable iff $\dim\mathrm{span}\{u_i\}\le r$.
In that case any $P^\star$ whose range contains $\mathrm{span}\{u_i\}$
achieves the global causal optimum
$S^{\mathrm{caus}}_i(P^\star)=\lambda_1(K_i)$ for every $i$.
\end{corollary}

\begin{proof}
Direct from Definition~\ref{def:cov} and Proposition~\ref{prop:localmax}.
\end{proof}

% =========================================================================
\section{The Fork: Coherence Implies Coverability Strictly}
\label{sec:fork}
% =========================================================================

The two structural conditions are not interchangeable.

\begin{theorem}[Coherence implies coverability strictly]\label{thm:fork}
For a single PSD family $\{T_t\}$ on $\R^d$:
\begin{enumerate}
\item[(i)] $r$-coherence implies $r$-coverability.
\item[(ii)] The converse fails in general.
\end{enumerate}
\end{theorem}

\begin{proof}
(i) Let $\{T_t\}$ be $r$-coherent with shared rank-$r$ subspace $V^\star$,
and fix $t$.
$V^\star$ attains the Ky--Fan maximum
$\tr(\Pi_{V^\star}T_t)=\sum_{i=1}^r\lambda_i(T_t)$.
Choose an orthonormal basis $\{v_1,\ldots,v_r\}$ of $V^\star$ that
diagonalises $\Pi_{V^\star}T_t\Pi_{V^\star}$ on $V^\star$; let
$\mu_k:=v_k^\top T_t v_k$ ordered $\mu_1\ge\cdots\ge\mu_r$, so
$\sum_{k=1}^r\mu_k=\sum_{i=1}^r\lambda_i(T_t)$.
Cauchy interlacing for the principal submatrix
$(v_k^\top T_t v_l)_{k,l}$ \cite{HornJohnson2013} gives
$\mu_k\le\lambda_k(T_t)$ for every $k$; combined with the sum equality
this forces $\mu_k=\lambda_k(T_t)$ for every $k$.
In particular $\mu_1=v_1^\top T_t v_1=\lambda_1(T_t)$, and the Rayleigh
quotient equality case gives $v_1\in\Emax(T_t)$.
Hence $V^\star\cap\Emax(T_t)\supseteq\mathrm{span}\{v_1\}\ne\{0\}$ for
every $t$, so $V^\star$ is an $r$-cover.

(ii) For the converse, take $d=3$, $r=2$ and
\begin{equation}\label{eq:forkwitness}
T_1=\mathrm{diag}(5,4,0),\qquad
T_2=\mathrm{diag}(5,0,4).
\end{equation}
Both have simple top eigenvalue $5$ with top eigenvector $e_1$, so any
rank-$2$ subspace containing $e_1$ is a $2$-cover; the family is
$2$-coverable.
However, the top-$2$ eigenspaces are
$\mathrm{Top}_2(T_1)=\mathrm{span}\{e_1,e_2\}$ and
$\mathrm{Top}_2(T_2)=\mathrm{span}\{e_1,e_3\}$, which are distinct.
Their intersection is $\mathrm{span}\{e_1\}$, of dimension $1<r=2$, so no
rank-$2$ subspace is contained in both top-$2$ eigenspaces; the family
is not $2$-coherent.
\end{proof}

\begin{remark}[Predictive--causal dissociation]
\label{rem:dissoc}
Theorem~\ref{thm:fork} establishes an algebraic gap between coherence and
coverability for a single PSD family.
In a setting where the predictive operators $\{M_t\}$ and causal operators
$\{K_t\}$ are different functions of the same task index (as is generally
the case: $M_t$ encodes a prediction objective and $K_t$ a local
intervention response), the two existence conditions decouple.
A common optimal projection under $S^{\mathrm{caus}}$ can exist while no
common optimal projection under $S^{\mathrm{pred}}$ does, and vice versa:
the existence law for a common optimal projection depends on the task
functional.
\end{remark}

% =========================================================================
\section{Weak Glue: Surrogate Optimization When Coverability Fails}
\label{sec:weakglue}
% =========================================================================

When the causal family $\{K_i\}$ has $\dim\mathrm{span}\{u_i\}>r$ (i.e.
fails $r$-coverability), no projector achieves global causal optimum
simultaneously.
The relevant question becomes \emph{approximate} causal projection
selection: how close to optimal can a single rank-$r$ projector come
across the family, and is there a tractable surrogate for the
non-Ky--Fan operator-norm objective?

\subsection{Setup and surrogate}

\emph{Standing assumption for this section.}
The regime is finite-support, $\mu=\sum_{i=1}^N w_i\delta_{x_i}$ with
$w_i>0$, $\sum w_i=1$, and each local Gram $A_i:=J_iJ_i^\top\succeq 0$
has \emph{simple} top eigenvalue ($\lambda_1(A_i)>\lambda_2(A_i)\ge 0$),
with top eigenvector $u_i$ and residual ratio
$\kappa_i:=\lambda_2(A_i)/\lambda_1(A_i)\in[0,1)$.
Repeated leading eigenspaces would require a projector-valued surrogate
in place of $u_i u_i^\top$ and are outside the scope of this section.

The capacity-first causal objective (the regime-averaged
operator-norm score) and its surrogate are:
\begin{align}
F(P) &:= \sum_{i=1}^N w_i\,\lambda_{\max}(P A_i P), \label{eq:F}\\
\tilde F(P) &:= \tr(P B_\lambda),\qquad
B_\lambda := \sum_{i=1}^N w_i\,\lambda_1(A_i)\,u_i u_i^\top.
\label{eq:Btilde}
\end{align}
The surrogate $\tilde F$ is a Ky Fan trace and is therefore maximised by
any rank-$r$ projection whose range is a dominant invariant subspace of
$B_\lambda$.
Let $P_F^\star\in\arg\max_{\rank P=r} F(P)$ (a true causal optimiser)
and $P_B^\star\in\arg\max_{\rank P=r}\tilde F(P)$ (a surrogate optimiser);
under strict $r$-gap of $B_\lambda$ both are unique up to range
equivalence on $\Gr(r,d)$, while under degeneracy both are interpreted as
\emph{any} element of the corresponding argmax.

\subsection{The sandwich and coarse regret}

\begin{theorem}[Weak-glue sandwich and regret]\label{thm:weakglue}
For all rank-$r$ projectors $P$:
\begin{equation}\label{eq:sandwich}
\tilde F(P)\;\le\;F(P)\;\le\;\tilde F(P)+\sum_{i=1}^N w_i\,\lambda_2(A_i).
\end{equation}
Consequently,
\begin{equation}\label{eq:coarse}
F(P_F^\star)-F(P_B^\star)\;\le\;\sum_{i=1}^N w_i\,\lambda_2(A_i).
\end{equation}
$\tilde F=F$ pointwise iff every $A_i$ is rank-$1$ (i.e.\ $\kappa_i=0$
for all $i$).
\end{theorem}

\begin{proof}
\emph{Per-term bound.}
Fix $i$ and let $A_i=\sum_{k=1}^d\lambda_k(A_i)e_k e_k^\top$ be the
spectral decomposition with $e_1=u_i$.
For any unit $v\in\ran(P)$,
\begin{equation*}
v^\top A_i v
=\sum_k\lambda_k(A_i)(e_k^\top v)^2
\le\lambda_1(A_i)(u_i^\top v)^2+\lambda_2(A_i)\sum_{k\ge 2}(e_k^\top v)^2
=\lambda_1(A_i)(u_i^\top v)^2+\lambda_2(A_i)\bigl(1-(u_i^\top v)^2\bigr).
\end{equation*}
Since $v\in\ran(P)$ and $P^\top=P$, $u_i^\top v=u_i^\top Pv=(Pu_i)^\top v$,
so by Cauchy--Schwarz $(u_i^\top v)^2\le\|Pu_i\|^2\|v\|^2=\|Pu_i\|^2$.
The map $t\mapsto\lambda_1(A_i)t+\lambda_2(A_i)(1-t)$ is non-decreasing on
$[0,1]$ since $\lambda_1\ge\lambda_2$, so
$v^\top A_iv\le\lambda_1(A_i)\|Pu_i\|^2+\lambda_2(A_i)(1-\|Pu_i\|^2)$.
Taking the supremum over unit $v\in\ran(P)$ gives the upper bound
\begin{equation}\label{eq:perterm}
\lambda_{\max}(PA_iP)\le\lambda_1(A_i)\|Pu_i\|^2+\lambda_2(A_i)(1-\|Pu_i\|^2).
\end{equation}
For the lower bound, when $\|Pu_i\|>0$ test $v=Pu_i/\|Pu_i\|\in\ran(P)$:
$v^\top A_i v=(Pu_i)^\top A_i(Pu_i)/\|Pu_i\|^2
\ge\lambda_1(A_i)(u_i^\top Pu_i)^2/\|Pu_i\|^2
=\lambda_1(A_i)\|Pu_i\|^2$,
using $u_i^\top Pu_i=\|Pu_i\|^2$ and dropping the nonnegative
$\lambda_k\ge 2$ terms.
When $\|Pu_i\|=0$ the bound is trivial.
Thus
\begin{equation}\label{eq:pertermlow}
\lambda_{\max}(PA_iP)\ge\lambda_1(A_i)\|Pu_i\|^2.
\end{equation}

\emph{Sandwich.}
Summing~\eqref{eq:pertermlow} weighted by $w_i$ gives $\tilde F(P)\le F(P)$;
summing~\eqref{eq:perterm} and using $1-\|Pu_i\|^2\le 1$ gives
$F(P)\le\tilde F(P)+\sum_i w_i\lambda_2(A_i)$.

\emph{Coarse regret~\eqref{eq:coarse}.}
Decompose
$F(P_F^\star)-F(P_B^\star)
=[F(P_F^\star)-\tilde F(P_F^\star)]+[\tilde F(P_F^\star)-\tilde F(P_B^\star)]
+[\tilde F(P_B^\star)-F(P_B^\star)]$.
The middle term is $\le 0$ by surrogate optimality of $P_B^\star$; the
third is $\le 0$ by the lower sandwich; the first is bounded by
$\sum_i w_i\lambda_2(A_i)$ from~\eqref{eq:perterm} summed.

\emph{Pointwise equality $\tilde F=F$.}
For each $i$, define the per-term gap
$\delta_i(P):=\lambda_{\max}(PA_iP)-\lambda_1(A_i)\|Pu_i\|^2\ge 0$
by~\eqref{eq:pertermlow}.
$F(P)-\tilde F(P)=\sum_i w_i\delta_i(P)$ with positive weights and
nonnegative summands, so $F=\tilde F$ pointwise iff $\delta_i(P)=0$ for
every $i$ and every rank-$r$ projector $P$.
Per $i$: take $P$ to project onto an eigenvector $e_k$ of $A_i$ with
$k\ne 1$.
Then $\lambda_{\max}(PA_iP)=\lambda_k(A_i)$ and $\|Pu_i\|^2=(e_k^\top u_i)^2=0$,
so $\delta_i(P)=\lambda_k(A_i)$.
Vanishing for all $k\ne 1$ forces $\lambda_k(A_i)=0$, i.e.\ $A_i$ is
rank-$1$.
Conversely, if every $A_i$ is rank-$1$, $A_i=\lambda_1(A_i)u_iu_i^\top$,
and $\lambda_{\max}(PA_iP)=\lambda_1(A_i)\|Pu_i\|^2$ exactly, so $F=\tilde F$.
\end{proof}

\subsection{Instance-dependent regret}

The coarse bound~\eqref{eq:coarse} treats every state's residual
spectral mass as fully adversarial.
A tighter bound accounts for the alignment of $P_F^\star$ with each
local peak direction $u_i$.

\begin{theorem}[Instance-dependent regret]\label{thm:instancereg}
\begin{equation}\label{eq:instreg}
F(P_F^\star)-F(P_B^\star)
\;\le\;
\sum_{i=1}^N w_i\,\lambda_2(A_i)\bigl(1-\|P_F^\star u_i\|^2\bigr).
\end{equation}
\end{theorem}

\begin{proof}
By the three-term decomposition in the proof of
Theorem~\ref{thm:weakglue}, $F(P_F^\star)-F(P_B^\star)\le
F(P_F^\star)-\tilde F(P_F^\star)$.
Applying the per-term upper bound~\eqref{eq:perterm} at $P=P_F^\star$ for
every $i$ and summing with weights $w_i$:
\[
F(P_F^\star)-\tilde F(P_F^\star)
=\sum_i w_i\bigl[\lambda_{\max}(P_F^\star A_i P_F^\star)-\lambda_1(A_i)\|P_F^\star u_i\|^2\bigr]
\le\sum_i w_i\lambda_2(A_i)\bigl(1-\|P_F^\star u_i\|^2\bigr).
\]
\end{proof}

\begin{remark}[Two mechanisms]
The instance-dependent bound factors regret at state $i$ as the product
of two independent mechanisms:
(a) residual spectral mass $\lambda_2(A_i)>0$, meaning $J_i$ carries
energy beyond its top singular direction;
(b) optimizer misalignment $1-\|P_F^\star u_i\|^2>0$, meaning the
true causal optimum fails to capture the local peak direction.
Either mechanism alone produces no regret: rank-$1$ Jacobians contribute
nothing regardless of alignment, and perfectly aligned states contribute
nothing regardless of $\lambda_2(A_i)$.
\end{remark}

\subsection{Equality conditions}

\begin{corollary}[Three exact regimes]\label{cor:equality}
$F(P_F^\star)-F(P_B^\star)=0$, i.e.\ the surrogate optimum equals the
true optimum, in each of the following cases:
\begin{enumerate}
\item[(R1)] \emph{Rank-$1$ Jacobians:} $\lambda_2(A_i)=0$ for every $i$.
Then $\tilde F=F$ pointwise.
\item[(R2)] \emph{Strong glue:} $\dim\mathrm{span}\{u_i\}\le r$, i.e.\
$\{A_i\}$ is $r$-coverable.
Both $P_F^\star$ and $P_B^\star$ achieve the global maximum
$\sum_i w_i\lambda_1(A_i)$.
\item[(R3)] \emph{Common-shape family:} $A_i=\alpha_i A$ for some fixed
$A\succeq 0$ and scalars $\alpha_i>0$.
Then $F$ and $\tilde F$ have identical maximiser sets, namely
$\{P:\rank(P)=r,\,\ran(P)\cap\Emax(A)\ne\{0\}\}$.
\end{enumerate}
\end{corollary}

\begin{proof}
(R1) is Theorem~\ref{thm:weakglue}.
(R2): $r$-coverability gives $\|P_B^\star u_i\|=1$ for every $i$ since
$P_B^\star$ contains $\mathrm{span}\{u_i\}$, so the lower sandwich gives
$F(P_B^\star)\ge\sum_i w_i\lambda_1(A_i)$.
The trivial upper bound is also $\sum_i w_i\lambda_1(A_i)$, so $P_B^\star$
achieves the global maximum.
(R3): $F(P)=(\sum_i w_i\alpha_i)\lambda_{\max}(PAP)$ and $B_\lambda
=(\sum_i w_i\alpha_i)\lambda_1(A)\,\Pi_{\Emax(A)}$, so both objectives are
proportional to the operator-norm criterion of the single matrix $A$
and have the same maximiser set
$\{P:\ran(P)\cap\Emax(A)\ne\{0\}\}$ by Proposition~\ref{prop:localmax}.
\end{proof}

In all three regimes the instance-dependent bound vanishes: either
$\lambda_2(A_i)=0$, or $\|P_F^\star u_i\|^2=1$ for every $i$.
Outside these regimes the surrogate optimum need not coincide with the
true optimum, but Theorem~\ref{thm:instancereg} controls the gap.

% =========================================================================
\section{Numerical Illustrations}
\label{sec:numerics}
% =========================================================================

This section gives three illustrations.
The first instantiates the predictive theorem in stationary LGDS.
The second is the fork witness from Theorem~\ref{thm:fork}.
The third is a $3\times 3$ weak-glue example showing the gap between the
coarse and instance-dependent regret bounds.
All numerical values are computed exactly: closed-form Bayes risk in the
LGDS case, exact eigendecomposition for the small-dimensional examples.
There is no Monte Carlo simulation.

\subsection{Predictive: LGDS instance}

\paragraph{Setup}
State dimension $d=8$, projection rank $r=2$.
$A=VDV^{-1}$ with eigenvalues drawn from $U[0.3,0.85]$ and random
non-orthogonal $V$ (seed $100$); $Q=0.1\,I_8$.
$\Sigma$ is computed from the discrete Lyapunov equation; $\Sigma\succ 0$
to machine precision.
Optimisation over $\Gr(2,8)$ uses $300$ random orthonormal
initialisations and a quasi-Newton iteration on $U$ followed by a thin
QR retraction to maintain $UU^\top=I_r$, with gradient-norm threshold
$10^{-12}$ and function-value tolerance $10^{-15}$.

\paragraph{Condition A: target-dependent}
Targets $T_1=([e_1^\top;e_2^\top],1)$ and $T_2=([e_5^\top;e_6^\top],5)$.
With non-normal $A$, $A^1$ and $A^5$ rotate state-space directions
differently; the principal angles between the top-$2$ eigenspaces of
$M_1$ and $M_2$ are $\CondAAngleOne^\circ$ and $\CondAAngleTwo^\circ$,
confirming that no single rank-$2$ projection is simultaneously
predictive-optimal.
By Corollary~\ref{cor:approxcoh}, the minimax predictive regret
$\Gamma_{\mathrm{pred}}>0$ certifies non-coherence; numerically,
$\Gamma_{\mathrm{pred}}=\GammaCondA>0$.
Table~\ref{tab:condA} gives the regret structure.

\begin{table}[tbp]
\caption{Predictive regret matrix, Condition A ($d=8$, $r=2$).
Values from the exact Bayes risk formula~\eqref{eq:lgdsbayes}.}
\label{tab:condA}
\begin{center}
\begin{tabular}{lcc}
\toprule
Projection & $\Reg$ on $T_1$ & $\Reg$ on $T_2$ \\
\midrule
$P^\star_1$ (optimal for $T_1$)         & \CondARegretUoneTa & \CondARegretUoneTb \\
$P^\star_2$ (optimal for $T_2$)         & \CondARegretUtwoTa & \CondARegretUtwoTb \\
$P_{\mathrm{joint}}$ (minimax)          & \CondARegretUjointTa & \CondARegretUjointTb \\
\bottomrule
\end{tabular}
\end{center}
\end{table}

\paragraph{Condition B: structurally coherent}
A target family is constructed by design to be $2$-coherent: normal
$A_{\mathrm{sym}}=VDV^\top$ with orthogonal $V$, common $C$ whose rows
are top-$2$ eigenvectors of $A$, varying horizon $\tau\in\{1,2,3,5\}$.
The four resulting $M_{C,\tau}$ share the same dominant
$2$-dimensional eigenspace by construction, and the four optimal
projections are identical: maximum pairwise principal angle is
$\CondBMaxAngleDeg^\circ$ to machine precision.
Under $\varepsilon=\CondBPertEps$ perturbations to $C$
($\CondBPertDraws$ random draws), the maximum principal angle is
$\CondBPertMaxRad$~rad, well below the $1.5\varepsilon=\CondBPertThreshold$~rad
threshold, illustrating that approximate coherence is quantified via
principal angles (Corollary~\ref{cor:approxcoh}).

\subsection{Fork: causal-without-predictive privilege}

The fork witness from~\eqref{eq:forkwitness}: $d=3$, $r=2$, $T_1
=\mathrm{diag}(5,4,0)$, $T_2=\mathrm{diag}(5,0,4)$.

\paragraph{Causal score (assigning $K_t=T_t$)}
Both $K_t$ have simple top eigenvalue $5$ with top eigenvector $e_1$.
Thus $\dim\mathrm{span}\{u_t\}=1\le 2=r$, so $\{K_t\}$ is
$2$-coverable.
Any rank-$2$ projection whose range contains $e_1$---e.g.
$P_{\mathrm{cov}}=\Pi_{\mathrm{span}\{e_1,v\}}$ for any unit
$v\perp e_1$---satisfies
$S^{\mathrm{caus}}_t(P_{\mathrm{cov}})=\lambda_1(K_t)=5$ for both $t$.
A common optimal rank-$2$ projection under $S^{\mathrm{caus}}$ exists.

\paragraph{Predictive score (assigning $M_t=T_t$)}
$\mathrm{Top}_2(M_1)=\mathrm{span}\{e_1,e_2\}$ and
$\mathrm{Top}_2(M_2)=\mathrm{span}\{e_1,e_3\}$.
The intersection is $\mathrm{span}\{e_1\}$, of dimension $1<r$.
Hence no rank-$2$ subspace is contained in both top-$2$ eigenspaces;
$\{M_t\}$ is not $2$-coherent, so no common optimal rank-$2$ projection
under $S^{\mathrm{pred}}$ exists.
The minimax predictive regret is
\begin{equation}
\Gamma_{\mathrm{pred}}^{\mathrm{fork}}
=\min_{P\in\Gr(2,3)}\max_{t\in\{1,2\}}
\bigl[\lambda_1(M_t)+\lambda_2(M_t)-\tr(P M_t)\bigr]
=\frac{1}{2}\cdot 4=2,
\end{equation}
attained by $P=\Pi_{\mathrm{span}\{e_1,(e_2+e_3)/\sqrt{2}\}}$, which
captures direction $e_1$ exactly and direction $(e_2+e_3)/\sqrt 2$
half-aligned with each of $e_2,e_3$.
The trace contribution from the second slot is $4\cdot\frac{1}{2}=2$ for
each target; the regret on each target is $4\cdot\frac{1}{2}=2$.
This dissociation is exact and structural.

\subsection{Weak-glue counterexample}

A two-state regime in $\R^3$ illustrates that the coarse regret bound
$\sum_i w_i\lambda_2(A_i)$ can grossly overestimate the actual regret,
while still detecting the qualitative point that the surrogate
optimiser need not equal the true optimiser.

\paragraph{Construction}
$N=2$, $d=3$, $r=1$, equal weights $w_1=w_2=\tfrac{1}{2}$, with
\begin{equation}\label{eq:wgex}
A_1=\mathrm{diag}(3,2,0),\qquad
A_2=\begin{pmatrix} 0.9 & 0.6 & 0\\ 0.6 & 0.9 & 0\\ 0 & 0 & 0\end{pmatrix}.
\end{equation}
$A_1$: $u_1=e_1$, $\lambda_1=3$, $\lambda_2=2$, $\kappa_1=2/3$.
$A_2$: $u_2=(e_1+e_2)/\sqrt 2$, $\lambda_1=1.5$, $\lambda_2=0.3$,
$\kappa_2=0.2$.

Restricting to rank-$1$ projectors $P_\theta=v(\theta)v(\theta)^\top$
with $v(\theta)=(\cos\theta,\sin\theta,0)$:
\begin{equation}
F(\theta)=\tfrac{1}{2}[3\cos^2\theta+2\sin^2\theta+0.9+0.6\sin(2\theta)]
=\tfrac{1}{2}[3.4+0.5\cos(2\theta)+0.6\sin(2\theta)].
\end{equation}
Setting $dF/d\theta=0$ gives $\tan(2\theta_F^\star)=1.2$, so
$\theta_F^\star\approx 25.1^\circ$ and
$F(P_F^\star)=\tfrac{1}{2}[3.4+\sqrt{0.5^2+0.6^2}]
=\tfrac{1}{2}(3.4+0.781)\approx 2.090$.

The surrogate matrix is
\begin{equation}
B_\lambda=\tfrac{1}{2}\Bigl[3\,e_1e_1^\top
+\tfrac{3}{2}\cdot\tfrac{(e_1+e_2)(e_1+e_2)^\top}{2}\Bigr]
=\begin{pmatrix} 1.875 & 0.375 & 0\\ 0.375 & 0.375 & 0\\ 0 & 0 & 0\end{pmatrix}.
\end{equation}
The top eigenvector of the leading $2\times 2$ block lies at
$\theta_B^\star\approx 13.3^\circ$ from $e_1$.
At $P_B^\star$:
\begin{equation}
F(P_B^\star)=\tfrac{1}{2}[3.4+0.5\cos 26.6^\circ+0.6\sin 26.6^\circ]
=\tfrac{1}{2}(3.4+0.447+0.268)\approx 2.058.
\end{equation}

\paragraph{Regret comparison}
The actual regret is
\begin{equation}
F(P_F^\star)-F(P_B^\star)\approx 2.090-2.058 = 0.032.
\end{equation}
The coarse bound from~\eqref{eq:coarse} is
\begin{equation}
\sum_i w_i\lambda_2(A_i)=\tfrac{1}{2}(2+0.3)=1.15,
\end{equation}
which exceeds the true regret by a factor of $1.15/0.032\approx 36$.
The instance-dependent bound from Theorem~\ref{thm:instancereg} is
computed using
$\|P_F^\star u_1\|^2=\cos^2(25.1^\circ)\approx 0.820$ and
$\|P_F^\star u_2\|^2=\bigl(\cos 25.1^\circ+\sin 25.1^\circ\bigr)^2/2
\approx 0.886$:
\begin{equation}
\sum_i w_i\lambda_2(A_i)(1-\|P_F^\star u_i\|^2)
\approx \tfrac{1}{2}\cdot 2\cdot 0.180+\tfrac{1}{2}\cdot 0.3\cdot 0.114
\approx 0.197.
\end{equation}
The instance-dependent bound is $0.197/0.032\approx 6$ times the actual
regret---substantially tighter than the coarse $36\times$, though
still not tight at $P_F^\star$.
The residual gap is due to the upper sandwich being itself loose
pointwise.
The picture: $P_B^\star\ne P_F^\star$ (separated by $\approx 11.8^\circ$),
yet the actual value gap is small because $F(\theta)$ has a broad
maximum.
The coarse bound cannot distinguish a direction error from a value
error; the instance-dependent bound partially can.
Table~\ref{tab:wg} summarises.

\begin{table}[tbp]
\caption{Weak-glue example: regret and bounds.
$N=2$, $d=3$, $r=1$, equal weights, $A_i$ as in~\eqref{eq:wgex}.}
\label{tab:wg}
\begin{center}
\begin{tabular}{lc}
\toprule
Quantity & Value \\
\midrule
$F(P_F^\star)$                                & $2.090$ \\
$F(P_B^\star)$                                & $2.058$ \\
Actual regret $F(P_F^\star)-F(P_B^\star)$     & $0.032$ \\
Coarse bound $\sum_i w_i\lambda_2(A_i)$       & $1.150$ \\
Instance-dependent bound, eq.~\eqref{eq:instreg} & $0.197$ \\
Coarse / actual                               & $\approx 36$ \\
Instance-dependent / actual                   & $\approx 6$ \\
$\theta_F^\star$, $\theta_B^\star$            & $25.1^\circ$, $13.3^\circ$ \\
\bottomrule
\end{tabular}
\end{center}
\end{table}

% =========================================================================
\section{Scope and Discussion}
\label{sec:scope}
% =========================================================================

\paragraph{What is earned}
For finite PSD families on $\R^d$ and a fixed projection rank $r$:
\begin{itemize}
\item Under the predictive trace score, a common optimal rank-$r$
projection exists iff the family is $r$-coherent
(Theorem~\ref{thm:pred}).
\item Under the causal operator-norm score, a common optimal rank-$r$
projection exists iff the family is $r$-coverable
(Theorem~\ref{thm:caus}).
\item $r$-coherence implies $r$-coverability strictly
(Theorem~\ref{thm:fork}).
\item When $\{A_i\}$ fails $r$-coverability, the eigenvalue-weighted
surrogate $B_\lambda$ converts $F$ into a Ky Fan eigenproblem, and
the surrogate regret is bounded by
$\sum_i w_i\lambda_2(A_i)(1-\|P_F^\star u_i\|^2)$, vanishing in three
explicit regimes (rank-$1$, strong glue, common-shape;
Corollary~\ref{cor:equality}).
\end{itemize}
The instance-dependent bound is structurally informative: it factors
regret into residual spectral mass and optimizer misalignment, and
attributes it pointwise to the states in the support of $\mu$.

\paragraph{What is not earned}
The results are stated for finite PSD families and finite-support
regimes.
Continuous regimes require replacing $B_\lambda$ with the integral
operator
$\E_{x\sim\mu}[\lambda_1(J_xJ_x^\top)\,u_xu_x^\top]$ and analysing it
spectrally; we have not proved a continuous-regime analogue.
The causal operators $K_t$ are taken as given PSD inputs; in nonlinear
settings $K_t$ would derive from a global intervention response rather
than a local Jacobian Gram, and the rank-$1$-driven local maximiser
characterisation (Proposition~\ref{prop:localmax}) need not extend.
The geometry-first pooled objective $G(P)=\lambda_{\max}(P\,\E[J_xJ_x^\top]P)$
is mentioned implicitly via the rank-$1$ regime where $G=F$; a full
treatment of the Jensen gap $F\ge G$ for $r>1$ is left to a companion
note.
Approximate coverability (a perturbation theory of $r$-coverability) is
open; the predictive analogue is
Corollary~\ref{cor:approxcoh}.

\paragraph{Practical implication}
Whenever a single rank-$r$ projection must serve a family of objectives,
the iff conditions of Theorems~\ref{thm:pred} and~\ref{thm:caus} are
testable: $r$-coherence by checking equality of top-$r$ eigenspaces;
$r$-coverability by computing $\dim\mathrm{span}\{u_i\}$ in the
simple-top case.
When neither exact condition holds, the two structural quantities
play different roles:
$\Gamma_{\mathrm{pred}}\ge 0$ is the minimax \emph{unavoidable
worst-case} predictive suboptimality of any single projection
(Corollary~\ref{cor:approxcoh}; $\Gamma_{\mathrm{pred}}>0$ certifies
non-coherence and quantifies the inherent task-family gap), while
the weak-glue expression
$\sum_i w_i\lambda_2(A_i)(1-\|P_F^\star u_i\|^2)$
\emph{upper-bounds} the loss incurred by replacing the true causal
optimiser $P_F^\star$ with the surrogate optimiser $P_B^\star$
(Theorem~\ref{thm:instancereg}).
$\Gamma_{\mathrm{pred}}=0$ is the certificate of structural
alignment between the system dynamics and the predictive task
family; $\Gamma_{\mathrm{pred}}>0$ is the certificate that dedicated
per-target projections are necessary for predictive optimality.

% =========================================================================
\section*{Acknowledgements}
% =========================================================================

The author received no external funding for this work.
Reproducible code for the LGDS predictive instance is available from a
publicly accessible repository\footnote{Repository details accompany the
non-blind version of the manuscript.}.
The fork witness and weak-glue counterexample are computed analytically
from the exact $3\times 3$ matrices specified in equations
\eqref{eq:forkwitness} and~\eqref{eq:wgex}.

\bibliographystyle{elsarticle-num}
\bibliography{refs}

\end{document}
